\section{Background and Related Work}
\label{sec:background}

\subsection{Reading Service Paths at the Network Layer}
\label{sec:background_cross_layer}

Users reach services through domain names, CDN resources, replicas, and
anycast targets. Routing selects the resulting path, and traceroute exposes
the intervening IP hops and AS transitions. These transitions are the usual
unit of path-diversity measurement: they are counted, compared across vantage
points, and used to describe how many distinct routes a service presents.

This is also the layer at which service reachability is commonly assessed.
Prior work measures the user's view of submarine criticality
~\cite{liu2020outofsight}, characterizes anycast deployment and its network
footprint~\cite{cicalese2015anycast,devries2017anycast}, and relates cable
infrastructure to AS-level structure~\cite{ye2023impact}. These measurements
describe network-layer reachability. Our question is whether diversity at that
layer corresponds to breadth across feasible physical directions.

The relationship is many-to-many. Different routers or AS transitions may
project onto one landing-region corridor, while a recurring network transition
may support several physical alternatives. A \emph{landing-region corridor}
is a direction-independent pair of bounded coastal regions connected by at
least one feasible cable candidate. Exact landing stations preserve facility
detail but split nearby access points serving the same direction; country-
level inventories merge coastlines and international directions. Landing
regions retain coastal structure between these two resolutions. Corridors are
the appropriate unit for comparing physical directions: exact cable identity
is often ambiguous, and counting several candidate cables through nearby
landings as independent directions would overstate the alternatives. The
aggregation preserves distinct coastal exit directions while consolidating
parallel candidate systems that represent the same direction.

\subsection{Mapping Paths onto Submarine Infrastructure}
\label{sec:background_existing_work}

Event-oriented studies connect new submarine infrastructure to observed
routing changes. Bischof et al. characterized Cuba's connectivity before and
after ALBA-1 and later formulated the task of connecting Internet observations
to the worldwide cable mesh~\cite{bischof2015cuba,bischof2018untangling}.
Fanou et al. measured routing changes following new cable deployments
~\cite{fanou2020unintended}. Liu et al. combine resource discovery, RIPE Atlas
traceroutes, geolocation, and propagation feasibility to quantify how users
reach Web resources through submarine paths~\cite{liu2020outofsight}.

Cross-layer cartography supplies the assignments used by our measurement.
Internet Atlas and InterTubes associate long-haul facilities with observed
paths~\cite{durairajan2013atlas,durairajan2015intertubes}. iGDB integrates
facilities, fiber, and network entities into a cross-layer map
~\cite{anderson2022igdb}. Nautilus associates IP links with candidate
submarine cables and attaches confidence to each assignment
~\cite{ramanathan2023nautilus}. Calypso combines cable layouts, network
relationships, latency, and inland connectivity to map routes and defines
Route Stress over those assignments~\cite{wang2025calypso}. Xaminer applies
cross-layer maps to disaster-oriented infrastructure analysis
~\cite{ramanathan2024xaminer}. Measurements of intercontinental links also
identify layer-3 links and gateway routers~\cite{carisimo2023hopaway}.

These systems produce an assignment from a link or path to physical
infrastructure that could plausibly carry it. Their validation evaluates the
assignment using cable failures, targeted traceroutes, operator maps, or
expert feedback. CLASP adopts this projection approach rather than claiming to
identify ground-truth cable use. The individual assignments are inputs to our
measurement; distributions aggregated over a path population are the objects
we compare.

\subsection{From Mapping to Distribution Measurement}
\label{sec:background_motivation}

An assignment does not by itself describe how observation mass is distributed
across physical directions or whether that distribution tracks the one visible
at the network layer. We separate two quantities that are easy to conflate.
\emph{Submarine exposure} is the fraction of valid traceroutes containing a
feasible inter-region corridor. \emph{Corridor concentration} describes how
candidate-bearing segment mass is distributed across corridors after entry.
A country can enter the candidate space frequently and distribute observations
widely, or enter infrequently with the exposed observations concentrated on
one direction.

Country and measurement identity jointly define each observation population.
Source geography and interconnection describe the context from which paths
start; anycast, replica selection, and routing determine the destinations and
paths observed during the measurement window. DNS Roots illustrate this
interaction because functionally equivalent services operate independent
anycast deployments. The three applications retain their own selected
destinations. The two dynamic multi-target measurements deliberately sample a
broader target population and provide topology-reference distributions against
which the service-facing distributions can be read.

CLASP aligns the statistical unit across layers. A network transition and its
corridor support originate from the same atomic segment, so their concentration
difference describes how a fixed observation population is organized. The
framework next defines this paired construction, and Section
~\ref{sec:datasets} records the aligned path and infrastructure inputs.